\subsection{Convenção da série e integrais auxiliares}
Para um sinal qualquer de período angular $L$, adota-se $k_n=2\pi n/L$. Dessa forma, sua representação em série de Fourier é dada por:
\begin{equation}
    \begin{aligned}
        v(x)&=a_0+\sum_{n=1}^{\infty}\bigl[
            a_n\cos(k_nx)+b_n\sin(k_nx)\bigr], \\
        a_0&=\frac{1}{L}\int_0^L v(x)\,dx, \\
        a_n&=\frac{2}{L}\int_0^L v(x)\cos(k_nx)\,dx, \\
        b_n&=\frac{2}{L}\int_0^L v(x)\sin(k_nx)\,dx.
    \end{aligned}
    \label{eq:convencao_fourier}
\end{equation}

Nessa representação, $a_0$ é a própria média.
Na reconstrução $v_{i,7}$, o índice 7 indica truncamento na ordem 7;
as ordens com coeficientes nulos permanecem incluídas. Não se trata
necessariamente de sete termos não nulos.

Para $n\geq1$, a amplitude de pico e a fase podem ser expressas por:
\begin{equation}
    \begin{aligned}
        C_n&=\sqrt{a_n^2+b_n^2}, \\
        \alpha_n&=\operatorname{atan2}(b_n,a_n), \\
        v_n(x)&=C_n\cos(k_nx-\alpha_n).
    \end{aligned}
\end{equation}

Se $C_n=0$, a fase é indeterminada e não altera o sinal. As amplitudes
$C_n$ são valores de pico; o valor eficaz do harmônico é $C_n/\sqrt{2}$.

As integrações por partes utilizam, para $k\neq0$:
\begin{equation}
    \begin{aligned}
        \int x\cos(kx)\,dx
            &=\frac{x\sin(kx)}{k}+\frac{\cos(kx)}{k^2}+C, \\
        \int x\sin(kx)\,dx
            &=-\frac{x\cos(kx)}{k}+\frac{\sin(kx)}{k^2}+C.
    \end{aligned}
\end{equation}

